\documentclass[11pt]{article}
\usepackage{amsmath,amssymb,amsthm,fullpage}
\usepackage[colorlinks=true,linkcolor=blue,citecolor=blue]{hyperref}

\newtheorem{theorem}{Theorem}
\newtheorem{proposition}[theorem]{Proposition}
\newtheorem{lemma}[theorem]{Lemma}
\newtheorem{corollary}[theorem]{Corollary}
\theoremstyle{definition}
\newtheorem{definition}[theorem]{Definition}
\newtheorem{remark}[theorem]{Remark}

\DeclareMathOperator{\Par}{Par}
\DeclareMathOperator{\CRST}{CRST}
\DeclareMathOperator{\CSSYT}{CSSYT}
\newcommand{\la}{\lambda}
\newcommand{\Z}{\mathbb{Z}}
\newcommand{\edot}{\dot{e}}

\title{No monomial statistic on cylindric tableaux carries Warnaar's\\
       generalised affine Jacobi--Trudi determinant}
\author{Clio}
\date{10 October 2026}

\begin{document}
\maketitle

\begin{abstract}
Warnaar [arXiv:2511.17034, \S6] asks whether Huh--Kim--Krattenthaler--Okada's
combinatorial evaluation of a $t=1$ affine bounded Littlewood determinant
[arXiv:2301.13117, Thm.~3.3] extends to his generalised determinant
\eqref{eq:gen} \emph{``by introducing an additional statistic on cylindric
tableaux''}. We prove that it does not, if ``statistic'' is read in the usual
sense of a $\Z$-valued function $T\mapsto\mathrm{stat}(T)$ contributing a single
power $t^{\mathrm{stat}(T)}$. The obstruction is a cardinality count, not a
sign or a symmetry. For every $\ell\ge0$, setting $M=\ell+2$ and
$k=1$, $n=2M$, the coefficient of $x_1\cdots x_n$ in Warnaar's determinant is
$C_M-(2M-1)t^{M-1}+t^{M+1}$ with $C_M$ a Catalan number, while the cylindric
tableaux available to carry it number $C_M$ with one sign and only $2M-2$ with
the other: the coefficient $-(2M-1)$ asks for one object more than exists.
At $\ell=0$ the failure is stronger still --- the coefficient is $2-3t+t^{3}$,
whose coefficients have absolute-value sum $6$, while the total number of
$(2,2)$-cylindric tableaux of that content, over \emph{all} shapes and not only
the two Huh--Kim--Krattenthaler--Okada use, is $4$ --- so not even re-choosing
the signs helps.
Consequently no twisted transfer-matrix trace $\mathrm{Tr}(D\,T^N)$ with a
diagonal twist $D$ can produce the determinant either: a twist contributes one
monomial per configuration, and the target needs more than one monomial's worth
per configuration. For $\ell\ge1$ the obstruction weakens in a way we make
precise: the strict form of the question (Thm.~3.3 with $t^{\mathrm{stat}(T)}$
inserted) still fails, but the slack created by the cylindric shapes that
Thm.~3.3 assigns weight~$0$ is enough to leave the re-signed form open.
Since no such slack exists at $\ell=0$, any uniform answer must attach a
\emph{polynomial} in $t$ to each cylindric tableau --- which is exactly the
shape of Macdonald's classical $\psi_T(t)$.
\end{abstract}

\section{Provenance}

The question addressed here is Warnaar's, quoted verbatim in \S\ref{sec:setup}.
The idea of looking for it on the integrable side --- as a diagonal twist of a
cylindric transfer matrix --- came to me from a conference talk; that talk is
\emph{not} an input to anything below and no construction here is attributed to
it. Every object, every definition and every step is taken from the two arXiv
sources and from Macdonald's book, all three of which I hold in full.
Neither Warnaar's paper nor Huh et al.'s contains the words ``transfer matrix''
or ``vertex model''; the integrable framing is mine, and in the end it is not
needed, because the obstruction proved below is upstream of it.

\section{Setup and the two verbatim inputs}\label{sec:setup}

Fix integers $k\ge1$, $\ell\ge0$, $n\ge1$ and set
\[
  w:=2\ell+2,\qquad M:=k+\ell+1,\qquad N:=2M=2k+w .
\]
Write $x=(x_1,\dots,x_n)$, let $e_r=e_r(x)$ be the elementary symmetric
functions ($e_r=0$ for $r<0$ or $r>n$), and let
$\edot_r:=e_r(x_1,x_1^{-1},\dots,x_n,x_n^{-1})$.
Following Huh--Kim--Krattenthaler--Okada (``HKKO'' below) put
\[
  f_r(x):=\sum_{i\in\Z}e_i(x)e_{i+r}(x),\qquad
  F_{r,N}(x):=\sum_{y\in\Z}f_{r+Ny}(x),
\]
so that $f_r=f_{-r}$ and $F_{r,N}=F_{-r,N}$.

\subsection*{Warnaar's object}

Warnaar \cite[\S6, eq.~(Eq\_generalise)]{W} asks for the expansions of
\begin{equation}\label{eq:gen}
  \mathcal G_{k,\ell,n}(x;t)\;:=\;
  \sum_{y\in\Z^{k}}\ \det_{1\le i,j\le k}
  \Bigl(t^{M y_i^{2}-j y_i}\bigl(
     \edot_{\,n-i+j-N y_i}(x)-\edot_{\,n+i+j-N y_i}(x)\bigr)\Bigr).
\end{equation}
He then writes, verbatim:
\begin{quote}\itshape
Recalling the notation \textup{(Eq\_FiN)}, \textup{(Eq\_generalise)} for $t=1$ is
\[
(x_1\cdots x_n)^{-k}\det_{1\le i,j\le k}\bigl(F_{i-j,2(k+\ell+1)}(x)-F_{i+j,2(k+\ell+1)}(x)\bigr),
\]
for which Huh et al.\ \textup{[HKKO25, Theorem 3.3]} obtained a combinatorial
expression in terms of cylindric tableaux as follows:
\[
(x_1\cdots x_n)^{-k}\Bigl(\ \sideset{}{'}\sum_{\la\in\Par_{n,2k}^{2\ell+2}}
-\sideset{}{''}\sum_{\la\in\Par_{n,2k}^{2\ell+2}}\Bigr)
\sum_{T\in\CSSYT_{n;2k,2\ell+2}(\la)}x^{T}.
\]
Here the prime in the first sum denotes the restriction that $\la$ must be even
and the double prime in the second sum denotes the restrictions that
$\la'_1-\la'_{2k}=2\ell+2$, the last $2\ell+2$ parts of $\la$ are odd and the
other parts are even.
\textbf{Is it possible to extend this result to \textup{(Eq\_generalise)} by
introducing an additional statistic on cylindric tableaux?}
\end{quote}
(Boldface mine. The question is the last sentence.)

\subsection*{HKKO Theorem 3.3}

HKKO \cite[Thm.~3.3]{HKKO}, in their own notation. For positive integers $h,w$
let $\Par(h,w)$ be the set of partitions $\la$ with $\ell(\la)\le h$ and
$\la_1-\la_h\le w$. For $\la$ with $\ell(\la)\le 2h$ and $\la_1-\la_{2h}\le w$ set
\begin{equation}\label{eq:cminus}
c^{\pm}_{2h,w}(\la)=
\begin{cases}
1,&\text{if }\la_{2i-1}=\la_{2i}\text{ for all }1\le i\le h,\\
\pm1,&\text{if }\la_1-\la_{2h}=w\text{ and }\la_{2i}=\la_{2i+1}\text{ for all }1\le i\le h-1,\\
0,&\text{otherwise.}
\end{cases}
\end{equation}
Then (their equation (3.6))
\begin{equation}\label{eq:C-}
  \sum_{\la\in\Par(2h,w)}c^{-}_{2h,w}(\la)\,s_{\la[2h,w]'}(x)
  \;=\;\det_{1\le i,j\le h}\bigl(F_{-i+j,\,2h+w}(x)-F_{i+j,\,2h+w}(x)\bigr),
\end{equation}
and, by \cite[Prop.~2.9]{HKKO},
$s_{\la[2h,w]'}(x)=\sum_{T\in\CRST(\la;2h,w)}x^{T}$, the sum being over
$(2h,w)$-cylindric row-strict tableaux of shape~$\la$. A tableau is
$(h,w)$-cylindric when it has at most $h$ rows and $T\cup(T+(h,-w))$ is again a
tableau of a valid skew shape; in particular $\CRST(\la;2h,w)=\emptyset$ unless
$\la\in\Par(2h,w)$, so \eqref{eq:C-} already ranges over \emph{all} cylindric
tableaux of this type. Throughout, $h=k$ and $N=2h+w=2k+2\ell+2$, which is
Warnaar's $2(k+\ell+1)$.

We also use HKKO's lattice-path form \cite[Prop.~2.7]{HKKO}: the number of
$T\in\CRST(\la;m,w)$ with content $\alpha$ equals the number of paths in $\Z^m$
from the origin to $(\la_1,\dots,\la_m)$ staying in
$\{x_1\ge x_2\ge\cdots\ge x_m\ge x_1-w\}$, the $i$-th step being a $0/1$ vector
with $\alpha_i$ ones.

\section{The gate, and that it is a gate}

\begin{lemma}\label{lem:t1}
$\edot_{\,n-r}(x)=(x_1\cdots x_n)^{-1}f_r(x)$ for all $r\in\Z$, and consequently
\[
 (x_1\cdots x_n)^{k}\,\mathcal G_{k,\ell,n}(x;1)
 =\det_{1\le i,j\le k}\bigl(F_{-i+j,N}-F_{i+j,N}\bigr)
 =\sum_{\la\in\Par(2k,w)}c^{-}_{2k,w}(\la)\!\!\sum_{T\in\CRST(\la;2k,w)}\!\!x^{T}.
\]
\end{lemma}

\begin{proof}
The first identity is Warnaar's \cite[proof of Lem.~5.2]{W}; it is also
immediate from $e_r(x^{\pm})=\sum_{m}e_m(x)\,e_{m+r-n}(x)\,(x_1\cdots x_n)^{-1}$.
For the second, the $(i,j)$ entry of the matrix in \eqref{eq:gen} at $t=1$ is,
by the first identity,
$(x_1\cdots x_n)^{-1}\bigl(f_{(i-j)+Ny_i}-f_{(i+j)-Ny_i}\bigr)$, in which the
dependence on $y$ is through $y_i$ alone. Hence the sum over $y\in\Z^{k}$ may be
carried into row~$i$ by multilinearity of the determinant in its rows, and
$\sum_{y\in\Z}f_{(i-j)+Ny}=F_{i-j,N}$,
$\sum_{y\in\Z}f_{(i+j)-Ny}=F_{i+j,N}$. Since $F_{r,N}=F_{-r,N}$ this is the
middle expression. The right-hand equality is \eqref{eq:C-} with $h=k$.
\end{proof}

\begin{remark}[where the difficulty is]
At $t=1$ the weight $t^{My_i^2-jy_i}$ is absent and the $y$-sum passes into a
single row. For $t\ne1$ that weight depends on the \emph{column} index $j$ as
well as on $y_i$, so it is not a row constant and the $y$-sum cannot be carried
in. This is precisely what Warnaar's question asks one to repair.
\end{remark}

Lemma~\ref{lem:t1} was verified symbolically (exact arithmetic, SymPy) at
$(k,\ell,n)=(1,0,4),(1,0,5),(1,1,5),(1,1,6),(1,1,8),(2,0,4),(2,0,5)$, and
\eqref{eq:C-} itself at $(h,w,n)=(1,1,4),(1,2,3),(1,2,4),(1,4,4),(1,4,5),(2,2,4)$.

\paragraph{The gate discriminates.} Three deliberately wrong weight functions
in place of $c^{-}_{2h,w}$ --- namely $c^{+}_{2h,w}$, the unsigned version, and
the version with the wall condition $\la_1-\la_{2h}=w$ replaced by $w-1$ ---
all \emph{fail} \eqref{eq:C-} at each of $(h,w,n)=(1,2,4),(1,4,5),(2,2,4)$,
while the true $c^{-}$ passes. So the $t=1$ comparison is not a constant
function of the claim.

\paragraph{The transcription of \eqref{eq:gen} is confirmed by a second
mechanism.} At $\ell=0$, \eqref{eq:gen} is Warnaar's Conjecture~1.5
(his (Eq\_PCn)), which asserts
$\mathcal G_{k,0,n}(x;t)=(x_1\cdots x_n)^{-k}\sum_{\la\text{ even},\ \la_1\le2k}P_\la(x;t)$.
Computing the right-hand side from Macdonald's symmetrisation formula for
Hall--Littlewood polynomials --- an engine sharing no code with
\eqref{eq:gen} --- confirms the identity exactly at $(n,k)=(3,1),(4,1),(3,2)$.
Six perturbations of the exponent $My_i^2-jy_i$ were tested against this gate:
replacing $-jy_i$ by $+jy_i$, deleting it, and changing $M$ to $M\pm1$ or
$My_i^2$ to $|y_i|$ all fail. One arm is \textbf{vacuous} and is recorded as
such: replacing $-jy_i$ by $-iy_i$ (column index by row index) \emph{passes} at
$k=1$ and $k=2$, so this gate cannot see which index carries the linear term.
Nothing below depends on that distinction.

\section{The obstruction}

\begin{definition}\label{def:stat}
Fix $(k,\ell,n)$ and let $\mathcal T$ denote the set of all $(2k,w)$-cylindric
row-strict tableaux with entries in $\{1,\dots,n\}$, of all shapes
$\la\in\Par(2k,w)$. For $T\in\mathcal T$ let $\la(T)$ be its shape.
\begin{itemize}
\item[(S)] Warnaar's question has a \emph{monomial solution} at $(k,\ell,n)$ if
there is $\mathrm{stat}:\mathcal T\to\Z$ with
\[
 (x_1\cdots x_n)^{k}\,\mathcal G_{k,\ell,n}(x;t)
 =\sum_{T\in\mathcal T}c^{-}_{2k,w}(\la(T))\;t^{\mathrm{stat}(T)}\,x^{T}.
\]
This is literally \eqref{eq:C-} with $t^{\mathrm{stat}(T)}$ inserted, i.e.\ the
display Warnaar quotes, graded.
\item[(S$^\pm$)] It has a \emph{signed monomial solution} if there are
$\varepsilon:\mathcal T\to\{0,\pm1\}$ and $\mathrm{stat}:\mathcal T\to\Z$ with
\[
 (x_1\cdots x_n)^{k}\,\mathcal G_{k,\ell,n}(x;t)
 =\sum_{T\in\mathcal T}\varepsilon(T)\,t^{\mathrm{stat}(T)}\,x^{T}.
\]
\end{itemize}
Clearly (S)$\Rightarrow$(S$^\pm$). Both put a single power of $t$ on each
cylindric tableau; (S$^\pm$) additionally allows the sign to be re-chosen
tableau by tableau and allows the shapes with $c^{-}=0$ to be recruited.
\end{definition}

For a content $\alpha=(\alpha_1,\dots,\alpha_n)\in\Z_{\ge0}^{n}$ write
\[
 c_\alpha(t):=[x^{\alpha}]\;(x_1\cdots x_n)^{k}\,\mathcal G_{k,\ell,n}(x;t)
 \in\Z[t,t^{-1}],\qquad
 \|c_\alpha\|_1:=\sum_{e}\bigl|[t^{e}]c_\alpha\bigr| ,
\]
and let
\[
 A_\alpha:=\#\{T\in\mathcal T:\ \mathrm{cont}(T)=\alpha,\ c^{-}(\la(T))=+1\},\quad
 B_\alpha:=\#\{\cdots=-1\},\quad
 \mathrm{TOT}_\alpha:=\#\{T\in\mathcal T:\mathrm{cont}(T)=\alpha\}.
\]
By Lemma~\ref{lem:t1}, $c_\alpha(1)=A_\alpha-B_\alpha$ for every $\alpha$.

\begin{lemma}[capacity]\label{lem:cap}
Fix $(k,\ell,n)$ and a content $\alpha$.
\begin{enumerate}
\item If \textup{(S)} holds then
 $\displaystyle\sum_{e}\max\bigl([t^{e}]c_\alpha,\,0\bigr)\le A_\alpha$ and
 $\displaystyle\sum_{e}\max\bigl(-[t^{e}]c_\alpha,\,0\bigr)\le B_\alpha$.
\item If \textup{(S$^\pm$)} holds then $\|c_\alpha\|_1\le \mathrm{TOT}_\alpha$.
\end{enumerate}
\end{lemma}

\begin{proof}
(1) Let $a_e$ (resp.\ $b_e$) be the number of $T\in\mathcal T$ of content
$\alpha$ with $c^-(\la(T))=+1$ (resp.\ $-1$) and $\mathrm{stat}(T)=e$. These are
nonnegative integers with $\sum_e a_e=A_\alpha$, $\sum_e b_e=B_\alpha$, and
comparing the coefficient of $x^{\alpha}t^{e}$ gives $[t^e]c_\alpha=a_e-b_e$.
Hence $\sum_e\max(a_e-b_e,0)\le\sum_e a_e=A_\alpha$, and symmetrically.
(2) Likewise $[t^e]c_\alpha=\sum_{T}\varepsilon(T)[\mathrm{stat}(T)=e]$ over the
$T$ of content~$\alpha$, so $\sum_e|[t^e]c_\alpha|\le\mathrm{TOT}_\alpha$.
\end{proof}

The content of Lemma~\ref{lem:cap} is that a monomial weight can only
\emph{redistribute} objects among powers of $t$; it cannot create them. The
$\ell^1$-norm of the answer is therefore bounded by the number of objects
available to carry it. We now exhibit a content where it is not.

\begin{theorem}[main]\label{thm:main}
Let $\ell\ge0$, put $M:=\ell+2$ and take
\[
  k=1,\qquad w=2\ell+2=2M-2,\qquad N=2M,\qquad n=N=2M,\qquad \alpha=(1^{n}).
\]
Then, with $C_M=\frac{1}{M+1}\binom{2M}{M}$ the $M$-th Catalan number,
\begin{equation}\label{eq:witness}
  c_\alpha(t)=C_M-(2M-1)\,t^{\,M-1}+t^{\,M+1},
  \qquad A_\alpha=C_M,\qquad B_\alpha=2M-2 .
\end{equation}
Consequently \textup{(S)} fails for \emph{every} $\ell\ge0$: Huh--Kim--Krattenthaler--Okada's
Theorem 3.3 with $t^{\mathrm{stat}(T)}$ inserted is false for every statistic
$\mathrm{stat}$, at every $\ell$.

Moreover, for $\ell=0$ (so $M=2$, $n=4$) one has
$\mathrm{TOT}_\alpha=A_\alpha+B_\alpha=4$ while $\|c_\alpha\|_1=6$, so
\textup{(S$^\pm$)} fails as well: no signed monomial statistic on cylindric
tableaux exists for Warnaar's Conjecture \textup{(Eq\_PCn)}.
\end{theorem}

\begin{proof}
\emph{Step 1: the coefficient.} For $k=1$ the determinant in \eqref{eq:gen} is
$1\times1$ with $i=j=1$, so by Lemma~\ref{lem:t1}'s first identity
\[
 (x_1\cdots x_n)\,\mathcal G_{1,\ell,n}(x;t)
 =\sum_{y\in\Z}t^{\,My^{2}-y}\bigl(f_{Ny}(x)-f_{Ny-2}(x)\bigr),
\]
where we used $f_{-r}=f_{r}$ to rewrite $\edot_{\,n+2-Ny}=(x_1\cdots x_n)^{-1}f_{Ny-2}$.
Since $n=N$ and $f_r=0$ for $|r|>n$, the term $f_{Ny}$ vanishes unless
$|y|\le1$, and $f_{Ny-2}$ vanishes unless $y\in\{0,1\}$ (for $y=-1$ one has
$|{-N-2}|=N+2>n$; for $|y|\ge2$, $|Ny-2|\ge 2N-2>N=n$ because $N\ge4$). The
surviving exponents are $My^2-y=0,\,M-1,\,M+1$ at $y=0,1,-1$, so
\begin{equation}\label{eq:threeterm}
 (x_1\cdots x_n)\,\mathcal G_{1,\ell,n}
 =\bigl(f_0-f_2\bigr)+t^{\,M-1}\bigl(f_{N}-f_{N-2}\bigr)+t^{\,M+1}f_{N}.
\end{equation}
Now extract $[x_1x_2\cdots x_n]$. For $a+b=n$ one has
$[x_1\cdots x_n]\,e_a e_b=\binom{n}{a}$, since the monomial is squarefree of
degree $n$ and the two factors must use complementary variable sets; and
$[x_1\cdots x_n]\,e_ae_b=0$ when $a+b\ne n$. Hence in
$f_r=\sum_m e_me_{m+r}$ exactly one term survives, the one with $2m+r=n$:
\[
 [x_1\cdots x_n]f_0=\binom{2M}{M},\quad
 [x_1\cdots x_n]f_2=\binom{2M}{M-1},\quad
 [x_1\cdots x_n]f_{N-2}=\binom{n}{1}=2M,\quad
 [x_1\cdots x_n]f_{N}=\binom{n}{0}=1 .
\]
(For $f_2$: $2m+2=n=2M$ gives $m=M-1$. For $f_{N-2}=f_{n-2}$: $m=1$. For
$f_N=f_n$: $m=0$.) Substituting into \eqref{eq:threeterm} and using the
standard identity $\binom{2M}{M}-\binom{2M}{M-1}=C_M$ gives
\[
 c_\alpha(t)=C_M+t^{\,M-1}\bigl(1-2M\bigr)+t^{\,M+1},
\]
which is \eqref{eq:witness}.

\emph{Step 2: the counts.} Content $\alpha=(1^n)$ forces $|\la|=n$, and by
HKKO's lattice-path description each of the $n$ steps is $(1,0)$ or $(0,1)$,
so the difference $d=x_1-x_2$ performs a $\pm1$ walk of length $n=2M$ from
$d=0$, confined to $[0,w]=[0,2M-2]$ by the region condition
$x_1\ge x_2\ge x_1-w$.

For $A_\alpha$ the endpoint is $d=0$ (shape $(M,M)$, the unique even shape of
size $n$ with $\la_1=\la_2$). A nonnegative $\pm1$ path of length $2M$
returning to $0$ is a Dyck path, and there are $C_M$ of them; its height is at
most $M\le 2M-2$ for $M\ge2$, so the ceiling is never active and no path is
lost. Hence $A_\alpha=C_M$.

For $B_\alpha$ the endpoint is $d=w=2M-2$ (shape $(2M-1,1)$). With $u$ up-steps
and $\delta$ down-steps, $u+\delta=2M$ and $u-\delta=2M-2$, so $\delta=1$: the
walk is all up except for a single down-step, say at position $p$. Its value
after $p$ steps is $p-2$, which must be $\ge0$, so $p\ge2$; and its maximum is
$\max(p-1,\,2M-2)$, which must be $\le 2M-2$, so $p\le 2M-1$. Both conditions
are also sufficient. Hence $p\in\{2,\dots,2M-1\}$ and $B_\alpha=2M-2$.

\emph{Step 3: (S) fails.} Suppose (S) held. In the notation of
Lemma~\ref{lem:cap}(1), comparing coefficients of $t^{M-1}$ gives
$a_{M-1}-b_{M-1}=-(2M-1)$, so $b_{M-1}\ge 2M-1$. But
$\sum_e b_e=B_\alpha=2M-2<2M-1$, a contradiction. This holds for every
$M\ge2$, i.e.\ every $\ell\ge0$.

\emph{Step 4: (S$^\pm$) fails at $\ell=0$.} Here $M=2$, $n=4$, $w=2$, and
\eqref{eq:witness} reads $c_\alpha(t)=2-3t+t^{3}$ with $A_\alpha=B_\alpha=2$.
The shapes of size $4$ in $\Par(2,2)$ are $(2,2)$ and $(3,1)$, since $(4,0)$
violates $\la_1-\la_2\le w=2$; these are exactly the shapes carrying
$c^{-}=+1$ and $c^{-}=-1$. Hence $\mathrm{TOT}_\alpha=A_\alpha+B_\alpha=4$,
while $\|c_\alpha\|_1=2+3+1=6>4$, and Lemma~\ref{lem:cap}(2) applies.
\end{proof}

\begin{remark}[why $\ell=0$ is the only place the stronger form bites]
The deficit in Step~3 is $(2M-1)-(2M-2)=1$, uniformly in $\ell$ --- the strict
form fails by exactly one tableau, always. The weaker form (S$^\pm$) is a
different matter: it needs $\|c_\alpha\|_1=C_M+2M$ to exceed
$\mathrm{TOT}_\alpha$, the number of \emph{all} confined walks of length $2M$
in $[0,2M-2]$. For $M=2$ the only reachable endpoints are $d=0$ and $d=2=w$,
both of which HKKO's theorem already uses, so $\mathrm{TOT}=A+B=4<6$. For
$M\ge3$ the intermediate even endpoints $d=2,4,\dots,w-2$ become available;
these are shapes with $c^{-}_{2,w}(\la)=0$, invisible to Theorem 3.3, and they
supply enough slack that $\|c_\alpha\|_1<\mathrm{TOT}_\alpha$ --- at $M=3$,
$\mathrm{TOT}=18$ against $\|c_\alpha\|_1=11$. So the escape route available
for $\ell\ge1$ is precisely ``recruit the shapes the theorem assigns weight
zero'', and at $\ell=0$ there are none to recruit in even degree.
\end{remark}

\begin{theorem}[the obstruction is not confined to $k=1$]\label{thm:k2}
At $(k,\ell,n)=(2,0,4)$ and content $\alpha=(2,2,2,2)$,
\[
 c_\alpha(t)=10-16t+11t^{2}-4t^{3},\qquad A_\alpha=6,\ B_\alpha=5,\
 \mathrm{TOT}_\alpha=11,\qquad \|c_\alpha\|_1=41 .
\]
Hence \textup{(S$^\pm$)} fails, by a margin of $30$. Across all $313$ monomials
at these parameters, $145$ violate \textup{(S$^\pm$)}.
\end{theorem}

\paragraph{Scope of the computation.} Theorem~\ref{thm:main} was found by, and is consistent with, an exhaustive
sweep. For each listed $(k,\ell,n)$ every
monomial of $(x_1\cdots x_n)^{k}\mathcal G$ was examined; the identity
$c_\alpha(1)=A_\alpha-B_\alpha$ was asserted for \emph{every} one of them (and
holds), and the two capacity inequalities were tested. See
Table~\ref{tab:sweep}. The counts $A_\alpha,B_\alpha,\mathrm{TOT}_\alpha$ were
computed twice by independent mechanisms --- from the cylindric Jacobi--Trudi
determinant \cite[Thm.~2.10]{HKKO} and from the lattice-path model
\cite[Prop.~2.7]{HKKO} --- agreeing on all $476$ contents compared at
$(k,w,n)=(1,2,4),(1,4,5),(2,2,4)$.

\begin{table}[h]\centering
\begin{tabular}{llrrrl}
\hline
$(k,\ell)$ & $n$ & monomials & (S) viol. & (S$^\pm$) viol. & worst witness $c_\alpha(t)$ \\
\hline
$(1,0)$ & $3$  & $14$    & $0$    & $0$    & --- \\
$(1,0)$ & $4$  & $41$    & $1$    & $1$    & $2-3t+t^{3}$; $A{=}B{=}2$, $\mathrm{TOT}{=}4$ \\
$(1,0)$ & $5$  & $122$   & $10$   & $10$   & $2-3t+t^{3}$; $\mathrm{TOT}{=}4$ \\
$(1,0)$ & $6$  & $365$   & $61$   & $61$   & $5-9t+5t^{3}-t^{6}$; $\mathrm{TOT}{=}8$ \\
$(1,1)$ & $5$  & $122$   & $0$    & $0$    & --- \\
$(1,1)$ & $6$  & $365$   & $1$    & $0$    & $5-5t^{2}+t^{4}$; $A{=}5,B{=}4,\mathrm{TOT}{=}18$ \\
$(1,1)$ & $7$  & $1094$  & $14$   & $0$    & $5-5t^{2}+t^{4}$ \\
$(1,1)$ & $8$  & $3281$  & $113$  & $0$    & $14-20t^{2}+7t^{4}$; $A{=}14,B{=}13$ \\
$(1,2)$ & $7$  & $1094$  & $0$    & $0$    & --- \\
$(1,2)$ & $8$  & $3281$  & $1$    & $0$    & $14-7t^{3}+t^{5}$; $A{=}14,B{=}6,\mathrm{TOT}{=}68$ \\
$(1,2)$ & $9$  & $9842$  & $18$   & $0$    & $14-7t^{3}+t^{5}$ \\
$(1,3)$ & $10$ & $29525$ & $1$    & $0$    & $42-9t^{4}+t^{6}$; $A{=}42,B{=}8,\mathrm{TOT}{=}250$ \\
$(2,0)$ & $4$  & $313$   & $145$  & $145$  & $10-16t+11t^{2}-4t^{3}$; $\mathrm{TOT}{=}11$ \\
$(2,0)$ & $5$  & $1563$  & $1051$ & $1051$ & $27-60t+55t^{2}-20t^{3}-5t^{4}+4t^{5}$; $\mathrm{TOT}{=}21$ \\
$(2,1)$ & $6$  & $7813$  & $1817$ & $0$    & $82-90t^{2}+46t^{3}-6t^{5}$; $A{=}76,B{=}44$ \\
\hline
\end{tabular}
\caption{Exhaustive per-monomial test of Lemma~\ref{lem:cap}. ``(S) viol.''
counts contents $\alpha$ admitting no $\mathrm{stat}$ with HKKO's signs;
``(S$^\pm$) viol.'' counts those admitting no signed monomial statistic at all.
For every monomial of every row the identity $c_\alpha(1)=A_\alpha-B_\alpha$ of
Lemma~\ref{lem:t1} was asserted, and holds. The rows with
$n=2\ell+4$ and worst witness at $\alpha=(1^{n})$ --- $(1,0)/4$, $(1,1)/6$,
$(1,2)/8$, $(1,3)/10$ --- are the instances of Theorem~\ref{thm:main}, whose
predicted witnesses $2-3t+t^{3}$, $5-5t^{2}+t^{4}$, $14-7t^{3}+t^{5}$,
$42-9t^{4}+t^{6}$ and counts $(C_M,\,2M-2)=(2,2),(5,4),(14,6),(42,8)$ all
appear. Note the dichotomy: at $\ell=0$ the two violation columns coincide; at
$\ell\ge1$ the second is empty.}
\label{tab:sweep}
\end{table}

\section{What this says about the twisted-trace ansatz}

The session's original target was to compute a twisted trace
$Z=\mathrm{Tr}(D\,T^{N})$ of a cylindric transfer matrix, $D$ a diagonal twist,
and ask whether the grading it induces on cylindric tableaux is Warnaar's
statistic. Theorem~\ref{thm:main} answers this without building the model, and for a
reason that is \emph{not} the reason I expected.

\begin{corollary}\label{cor:twist}
Suppose a vertex model has configurations in bijection with the
$(2k,w)$-cylindric tableaux, with Boltzmann weight $x^{T}$, and let $D$ be any
operator acting on the winding sector of a configuration by a scalar
$t^{\,\varphi(\cdot)}$. Then $\mathrm{Tr}(D\,T^{N})$ has the form
$\sum_{T}\varepsilon(T)t^{\mathrm{stat}(T)}x^{T}$ and is therefore subject to
Lemma~\ref{lem:cap}(2). At $(k,\ell,n)=(1,0,4)$ it cannot equal
$(x_1\cdots x_n)^{k}\mathcal G$.
\end{corollary}

That is: the twist ansatz dies, but not because the twist is \emph{diagonal},
and not because its exponent is linear while Warnaar's is quadratic. It dies on
\emph{cardinality}. A twist --- diagonal or not, linear or quadratic in the
winding --- contributes one monomial in $t$ per configuration, and Warnaar's
determinant needs more than one monomial's worth per configuration. The
discrepancy named by the brief turns out to be upstream of the integrable
structure entirely.

The secondary, weaker observation remains true and is worth recording, because
it says where the quadratic exponent comes from.

\begin{proposition}\label{prop:lin}
Let $V=\bigoplus_{a\in L}V_a$ be graded by an abelian group $L$, let $T$ be any
operator on $V$, and let $D$ act on $V_a$ by $t^{\langle\varphi,a\rangle}$ for a
homomorphism $\varphi\colon L\to\Z$. Then
$\mathrm{Tr}(D\,T^{N})=\sum_{a\in L}t^{\langle\varphi,a\rangle}\mathrm{Tr}_{V_a}(\Pi_a T^{N})$,
so the exponent is additive in the sector label. Warnaar's exponent
$Q(y)=M\sum_i y_i^{2}-\sum_i jy_i$ is not: $Q(2e_1)-2Q(e_1)=2M\ne0$. Hence if
the lattice variable $y$ of \eqref{eq:gen} is identified with the winding sector
of a cylindric transfer matrix --- the only natural identification, $y$ being
the translation part of the affine Weyl group acting on the cylinder --- then no
diagonal twist reproduces it. A positive-definite quadratic form on the winding
lattice is the signature of a level-$M$ theta weight, i.e.\ of an energy grading
$t^{L_0}$ with $L_0\supset\frac{1}{4M}J_0^{2}$, not of a flavour twist
$t^{J_0}$.
\end{proposition}

I flag the hypothesis explicitly: Proposition~\ref{prop:lin} is conditional on
identifying $y$ with the winding sector, which I have not proved. It is not
needed for Corollary~\ref{cor:twist}, which is unconditional.

\section{What the statistic must be instead}

Theorem~\ref{thm:main} says something sharper than ``no''.

At $\ell\ge1$ the strict reading of Warnaar's question --- insert
$t^{\mathrm{stat}(T)}$ into the display he quotes --- is false, but only just:
the deficit is exactly one tableau, at every $\ell$. And for $\ell\ge1$ the
re-signed form survives the capacity test with room to spare, the room coming \emph{entirely} from the cylindric
shapes $\la\in\Par(2k,w)$ with $c^{-}_{2k,w}(\la)=0$ --- at $(1,1,6)$, the shape
$(4,2)$, which contributes $9$ of the $18$ tableaux. So for $\ell\ge1$ a
monomial statistic is not excluded, provided one is willing to recruit shapes
that HKKO's theorem does not mention and to choose signs that are not
$c^{-}$.

At $\ell=0$ that escape is closed: there are no intermediate shapes in even
degree. So any answer \emph{uniform in $\ell$} --- and Warnaar's question is
posed uniformly, $\ell$ being a free nonnegative integer --- must abandon the
monomial form. It must attach to each cylindric tableau a Laurent
\emph{polynomial} $\pi_T(t)\in\Z[t^{\pm1}]$, not a single power.

This is not an exotic requirement. It is the shape of the classical answer one
level down. Macdonald \cite[III, (5.8$'$), (5.9$'$), (5.11$'$), pp.~229--230]{M}
gives
\[
 P_{\la/\mu}(x;t)=\sum_{T}\psi_T(t)\,x^{T},\qquad
 \psi_T(t)=\prod_{i\ge1}\psi_{\la^{(i)}/\la^{(i-1)}}(t),\qquad
 \psi_{\la/\mu}(t)=\prod_{j\in J}\bigl(1-t^{m_j(\mu)}\bigr),
\]
summed over tableaux of shape $\la/\mu$ --- a product of factors $1-t^{m}$ per
tableau, a polynomial, not a power of $t$. At $\ell=0$ Warnaar's determinant
\emph{is} (conjecturally, by his (Eq\_PCn), which I verified at
$(n,k)=(3,1),(4,1),(3,2)$) a sum of Hall--Littlewood polynomials
$\sum_{\la\text{ even},\la_1\le2k}P_\la(x;t)$. So the $t$ in the problem is a
Hall--Littlewood bulk parameter, which deforms a tableau count by a product of
$(1-t^{m})$'s; it is not a boundary fugacity, which would deform it by a power.
The quadratic Gaussian $t^{My_i^{2}-jy_i}$ on the Weyl-sum side of the identity
is what such a bulk weight \emph{sums to} after the affine folding, not what it
\emph{is} on the tableau side.

\section{Gaps and what is not claimed}

\begin{enumerate}
\item Theorem~\ref{thm:main} is a complete proof of a negative statement, uniform in
$\ell$, with every number derived in the proof; Theorem~\ref{thm:k2} is a
single verified parameter point. They rest on exactly two external inputs:
HKKO Theorem~3.3 (eq.~\eqref{eq:C-}) and the definition \eqref{eq:gen}. They do
\emph{not} rest on Warnaar's Conjecture (Eq\_PCn), which is used only as an
independent check on my transcription.
\item I have \emph{not} shown that (S$^\pm$) holds for any $\ell\ge1$. The
capacity test is necessary, not sufficient; failing to violate it is not
evidence of existence. Table~\ref{tab:sweep} records a non-violation, which is
an absence of evidence.
\item I have \emph{not} constructed the polynomial weight $\pi_T(t)$ that
\S6 argues must exist. Identifying it is the real form of Warnaar's question
and is open.
\item Proposition~\ref{prop:lin} is conditional, as flagged there.
\item Trap~5 of the session brief (Orevkov--Orevkov's $n\le14$ degeneracy for
multiplicative versus additive polytopes) is recorded as \textbf{not
applicable}: nothing here is a honeycomb or polytope computation.
\end{enumerate}

\begin{thebibliography}{9}
\bibitem{HKKO} J.~Huh, J.~S.~Kim, C.~Krattenthaler, S.~Okada,
\emph{Bounded Littlewood identities for cylindric Schur functions},
arXiv:2301.13117. Theorem~3.3 and \S2 as cited above.
\bibitem{M} I.~G.~Macdonald, \emph{Symmetric Functions and Hall Polynomials},
2nd ed., Oxford University Press. Chapter III, \S5, pp.~229--230.
\bibitem{W} S.~O.~Warnaar, arXiv:2511.17034. \S6 (Open problems), and
Conjecture (Eq\_PCn).
\end{thebibliography}

\end{document}
